%% ============================================================
%%  Chapitre 3 — Moteur de Signaux (Signal Engine A→G)
%% ============================================================
\chapter{Moteur de Signaux Personnalisé}
\label{ch:signaux}

\section{Philosophie~: Un Signal n'est Pas une Stratégie}
\label{sec:philosophie_signal}

\subsection{La distinction fondamentale}

Cette distinction n'est pas sémantique~: elle est le fondement architectural de
l'ensemble du moteur de signaux. López de Prado~\cite{lopezdeprado2018} la formule
sans ambiguïté~:

\begin{quote}
\textit{«~[Feature analysts] collect and catalogue libraries of findings that can be
useful to a multiplicity of stations. [...] A common error is to believe that feature
analysts develop strategies. Instead, feature analysts collect and catalogue libraries
of findings.~»} — López de Prado (2018), Ch.~1, p.~7.
\end{quote}

Le tableau~\ref{tab:signal_vs_strategie} synthétise cette distinction.

\begin{table}[H]
\centering
\caption{Signal versus Stratégie~: distinction architecturale fondamentale}
\label{tab:signal_vs_strategie}
\begin{tabularx}{\linewidth}{@{}lYY@{}}
\toprule
                    & \textbf{Signal}                    & \textbf{Stratégie}                       \\
\midrule
\textbf{Définition} & Relation mesurée entre un indicateur et les retours futurs & Système complet agissant sur des signaux \\
\textbf{Sortie}     & Intention directionnelle~: $+1$, $-1$, $0$ & Ordres de trading avec sizing et risque \\
\textbf{Portée}     & Un indicateur, un paramètre, un horizon & Signaux multiples, construction de portefeuille \\
\textbf{Évaluation} & Sharpe de sous-période, stabilité, drawdown par fenêtre & PnL, retours ajustés du risque, métriques opérationnelles \\
\textbf{Analogie}   & La lecture du thermomètre & La décision d'allumer le chauffage \\
\bottomrule
\end{tabularx}
\end{table}

Le moteur de signaux produit des objets \texttt{VariantCurrentSignal} — des signaux
évalués, filtrés et pondérés par robustesse. Il ne produit \emph{pas} d'ordres de
trading, de tailles de position ou de poids de portefeuille. Ces responsabilités
sont en aval~: le Dashboard (Chapitre~\ref{ch:dashboard}) et, pour la validation statistique de l'edge, la page Analytics (Chapitre~\ref{ch:analytics}).

\subsection{Les cinq piliers initiaux}

Lors de la conceptualisation initiale du projet (Semaine~1~\cite{rapport_s1}), les
indicateurs ont été organisés en cinq piliers complémentaires~:

\begin{enumerate}
    \item \textbf{Tendance / Structure}~: Ichimoku, Moyennes Mobiles (SMA/EMA),
    Donchian Channels~;
    \item \textbf{Momentum / Timing}~: RSI, Stochastique, MACD, ADX~;
    \item \textbf{Volatilité / Risque}~: ATR, Bollinger, Keltner~;
    \item \textbf{Volume / Flux}~: OBV, MFI, A/D, Order Flow~;
    \item \textbf{Niveaux / Exécution}~: VWAP, Fibonacci, Parabolic SAR.
\end{enumerate}

\textbf{Règle d'or}~: complémentarité entre piliers, pas de \og vote brut \fg{},
puis confluence pour la décision~\cite{rapport_s1}.

\medskip

Cette vision initiale a évolué vers une architecture plus rigoureuse de 20 familles
évaluées par validation de stabilité temporelle, présentée dans la section suivante.

%% ============================================================
\section{Le Problème des Tests Multiples}
\label{sec:tests_multiples}

\subsection{Formalisation}

Si l'on évalue $k$ configurations d'indicateurs indépendantes sur un même historique,
sous l'hypothèse nulle ($H_0$~: l'indicateur n'a aucun pouvoir prédictif), le nombre
espéré de faux positifs au seuil $\alpha$ est donné par l'équation~\eqref{eq:faux_pos_ch2}~:

\begin{equation}
E\left[\text{Faux Positifs}\right] = \alpha \cdot k
\label{eq:faux_pos_ch2}
\end{equation}

Dans notre contexte~: 20 familles $\times$ 30 variantes $\times$ 3 horizons $= 1\,800$
évaluations potentielles. Avec $\alpha = 5\,\%$, on attendrait \textbf{90 faux positifs}
si aucune défense n'était mise en place.

\subsection{Le fondement information-théorique}

Un signal $s(t)$ a du pouvoir prédictif pour les retours futurs $y(t+h)$ si et
seulement si leur information mutuelle est strictement positive~:

\begin{equation}
I\bigl(s(t)\,;\,y(t+h)\bigr) > 0
\label{eq:info_mutuelle}
\end{equation}

En pratique, l'information mutuelle n'est pas directement calculable. Le pipeline
A→G utilise des proxys empiriques~: Sharpe de sous-période, stabilité, cohérence et drawdown
(Sections~\ref{sec:layer_c}).

\subsection{Le coût de friction}

Même un signal à pouvoir prédictif positif peut être \emph{non profitable} après
frais de transaction. L'edge minimal requis pour couvrir un aller-retour à un
coût de friction de $c$ (en fraction du prix) est~:

\begin{equation}
E\bigl[r_{t+h}\bigr] > 2c
\label{eq:edge_minimal}
\end{equation}

Ce coût de friction est intégré dans la validation par sous-périodes du Layer~B (Section~\ref{sec:layer_b}).

%% ============================================================
\section{Vue d'Ensemble~: Le Pipeline en 7 Couches}
\label{sec:pipeline_vue}

\begin{figure}[H]
\centering
\includegraphics[width=0.95\linewidth]{generated/ch3_pipeline_ag.png}
\caption{Le pipeline Signal Engine en 7 couches (Layers A→G).
Source~: documentation locale du moteur de signaux~\cite{repo_signal_docs}.}
\label{fig:pipeline_ag}
\end{figure}

Le tableau~\ref{tab:pipeline_defenses} synthétise, pour chaque couche, le mécanisme
de défense contre l'overfitting intégré.

\begin{table}[H]
\centering
\caption{Défenses anti-overfitting par couche du pipeline}
\label{tab:pipeline_defenses}
\begin{tabularx}{\linewidth}{@{}lYY@{}}
\toprule
\textbf{Layer} & \textbf{Mécanisme} & \textbf{Effet} \\
\midrule
A — Candidats       & Grille structurée, motivée économiquement & Évite la recherche aléatoire dans un espace infini \\
B — Sous-Périodes   & Fenêtres de test indépendantes, pas d'optimisation & Validation de stabilité temporelle (non OOS Pardo) \\
C — Robustesse      & Score multi-composante (4 métriques)        & Résiste au Goodhart's Law sur une métrique unique \\
D — Survivants      & Gate absolu 0.25 + percentile 60\%          & Élimine le \og best of a bad lot \fg{} \\
E — Redondance      & Clustering corrélation Pearson seuil 0.85   & Évite de compter la même information deux fois \\
G — Ensemble        & Moyenne pondérée multi-représentants        & Réduit la variance, dilue le bruit \\
\bottomrule
\end{tabularx}
\end{table}

%% ============================================================
\section{Layer A --- L'Univers des Candidats}
\label{sec:layer_a}

\subsection{Pourquoi une grille structurée~?}

Chaque candidat supplémentaire augmente le coût des tests multiples (équation~\eqref{eq:faux_pos_ch2}).
La grille est donc délibérément \emph{petite et structurée}~: exactement
\textbf{30 candidats} par couple (famille, horizon), chaque point de paramètre
ayant une justification économique explicite.

\subsection{Les 20 familles d'indicateurs}

Les 20 familles sont organisées en 4 catégories, chacune répondant à une question
de marché distincte~:

\begin{table}[H]
\centering
\caption{Les 20 familles d'indicateurs du Signal Engine}
\label{tab:20familles}
\begin{tabularx}{\linewidth}{@{}llYl@{}}
\toprule
\textbf{Catégorie} & \textbf{Famille} & \textbf{Règle signal ($+1$ = achat)} & \textbf{Type} \\
\midrule
\multirow{5}{*}{\textbf{Tendance}}
 & SMA               & $\text{Close} > \text{SMA}(w)$                         & État \\
 & EMA               & $\text{Close} > \text{EMA}(w)$                         & État \\
 & EMA Cross         & $\text{EMA}(f) > \text{EMA}(s)$                        & État \\
 & Ichimoku          & $\text{Close} > \text{Nuage}$ et Tenkan $>$ Kijun       & État \\
 & PSAR              & $\text{Close} > \text{SAR}$                             & État \\
\midrule
\multirow{5}{*}{\textbf{Momentum}}
 & MACD              & MACD Line $>$ Signal Line                               & Événement \\
 & ROC               & $\text{ROC}(n) > 0$                                     & État \\
 & TRIX              & $\text{TRIX}(n) > 0$                                    & État \\
 & ADX               & $+\text{DI} > -\text{DI}$ et $\text{ADX} > \text{seuil}$ & État \\
 & TSI               & $\text{TSI} > 0$                                        & État \\
\midrule
\multirow{5}{*}{\textbf{Oscillation}}
 & RSI               & $\text{RSI} < \text{survendu}$                          & Événement \\
 & Stochastique      & $\%K < 20$ et $\%K > \%D$                               & Événement \\
 & CCI               & $\text{CCI} < -100$                                     & Événement \\
 & MFI               & $\text{MFI} < \text{survendu}$                          & Événement \\
 & UO                & $\text{UO} < 30$                                        & Événement \\
\midrule
\multirow{5}{*}{\textbf{Volume}}
 & OBV               & $\text{OBV} > \text{SMA}(\text{OBV}, w)$                & État \\
 & CMF               & $\text{CMF}(w) > 0$                                     & État \\
 & A/D               & A/D Line en hausse                                      & État \\
 & VWAP              & $\text{Close} > \text{VWAP}(w)$                         & État \\
 & Force Index       & $\text{FI}(w) > 0$                                      & État \\
\bottomrule
\end{tabularx}
\end{table}

\subsection{Les horizons temporels}

Trois horizons sont définis, chacun correspondant à une durée de thèse économique~:

\begin{itemize}
    \item \textbf{Court terme}~: 1 semaine à 1 mois (5--20 barres journalières)~;
    \item \textbf{Moyen terme}~: 1 mois à 3--4 mois (21--80 barres)~;
    \item \textbf{Long terme}~: 6 mois à 1 an (120--250 barres).
\end{itemize}

\begin{table}[H]
\centering
\caption{Grille de paramètres de la famille SMA par horizon (30 valeurs chacun)}
\label{tab:sma_grid}
\begin{tabularx}{\linewidth}{@{}lY@{}}
\toprule
\textbf{Horizon} & \textbf{Fenêtres $w$ (30 valeurs)} \\
\midrule
Court terme  & 3, 5, 7, 8, 10, 12, 13, 15, 17, 18, 20, 22, 23, 25, 27, 28, 30, 33, 35, 38, 40, 42, 45, 48, 50, 55, 60, 65, 75, 90 \\
Moyen terme  & 10, 15, 18, 20, 22, 25, 27, 30, 33, 35, 38, 40, 45, 48, 50, 55, 60, 65, 70, 75, 80, 90, 100, 110, 120, 130, 150, 175, 200, 250 \\
Long terme   & 20, 30, 40, 50, 55, 60, 65, 70, 75, 80, 90, 100, 110, 120, 130, 140, 150, 160, 175, 190, 200, 220, 240, 260, 280, 300, 330, 350, 375, 400 \\
\bottomrule
\end{tabularx}
\end{table}

\subsection{Implémentation~: pattern registre décorateur}

Les familles sont enregistrées via un pattern décorateur, garantissant
l'extensibilité~:

\begin{lstlisting}[caption={Pattern d'enregistrement des familles d'indicateurs (Layer A)},label={lst:register_family}]
@register_family(family="sma")
def _sma_candidates(horizon: str) -> list[VariantDef]:
    windows = _SMA_WINDOWS[horizon]
    return [
        VariantDef(
            variant_id=f"sma_w{w}",
            family="sma",
            archetype="price_vs_sma",
            params={"window": w},
            horizon=horizon,
            description=f"Prix vs SMA({w})"
        )
        for w in windows
    ]
\end{lstlisting}

%% ============================================================
\section{Layer B --- Validation par Sous-Périodes}
\label{sec:layer_b}

\subsection{Le contrat signal}

Chaque fonction de signal doit respecter un contrat strict, vérifié après chaque
calcul~:

\begin{enumerate}
    \item \textbf{Longueur}~: $|\text{sortie}| = |\text{close}|$ — un signal par barre.
    \item \textbf{Valeurs}~: chaque élément $\in \{-1.0,\; 0.0,\; +1.0\}$ — pas de
    positions fractionnaires.
    \item \textbf{Warmup}~: les premières $w_{\min}$ barres sont remplies de $0.0$
    (position plate pendant la période de chauffe de l'indicateur).
    \item \textbf{Gestion NaN}~: tout NaN en sortie est remplacé par $0.0$.
\end{enumerate}

\subsection{Les fenêtres de sous-périodes indépendantes}

\begin{figure}[H]
\centering
\includegraphics[width=0.95\linewidth]{generated/ch3_subperiod_windows.png}
\caption{Fenêtres de sous-périodes pour l'évaluation de la stabilité (Layer B).
Chaque variante est évaluée avec ses paramètres fixes sur $n$ fenêtres indépendantes.}
\label{fig:oos_windows}
\end{figure}

Il est important de noter que ce mécanisme \textbf{n'utilise pas un Out-Of-Sample (OOS) traditionnel} au sens de Pardo. Les paramètres de la grille (Layer A) sont fixes et ne subissent \textbf{aucune optimisation}. L'historique complet est découpé en sous-périodes consécutives (fenêtres de test), et le "train window" n'est utilisé que comme un "burn-in" ou préchauffage de l'indicateur. Il s'agit d'une validation de stabilité temporelle.

\subsection{Métriques calculées par sous-période}

Pour chaque variante $v$ et chaque sous-période $k$, on calcule~:

\begin{itemize}
    \item $\text{Sharpe}_k(v)$~: ratio de Sharpe annualisé sur la sous-période~;
    \item $\text{MDD}_k(v)$~: Maximum Drawdown sur la sous-période~;
    \item $\mathbb{1}[\text{Sharpe}_k(v) > 0]$~: indicateur binaire de sous-période
    rentable.
\end{itemize}

Ces métriques alimentent directement le Layer~C.

%% ============================================================
\section{Layer C --- Score de Robustesse Multi-composante}
\label{sec:layer_c}

\subsection{Le gate de viabilité}

Avant le calcul du score, chaque variante doit passer une porte de viabilité~:

\begin{equation}
\text{viable} =
\bigl(n_{\text{fenêtres valides}} \geq 3\bigr)
\;\wedge\;
\left(\frac{\#\{k : \text{Sharpe}_k > 0\}}{n_{\text{fenêtres}}} \geq 0.40\right)
\label{eq:viability_gate}
\end{equation}

\textbf{Interprétation~:}
\begin{itemize}
    \item $\geq 3$ fenêtres valides~: en dessous, la moyenne de Sharpe n'a pas de
    signification statistique~;
    \item $\geq 40\,\%$ de fenêtres positives~: un signal qui fonctionne dans moins
    de 40\,\% des périodes est probablement du bruit.
\end{itemize}

Si la porte échoue~: $\text{reliability\_score} = 0.0$ et la variante est exclue
de toutes les couches suivantes.

\subsection{Le score de fiabilité composite}

Le score de fiabilité combine quatre composantes, chacune normalisée dans $[0, 1]$~:

\begin{equation}
\text{reliability} = 0.35 \cdot S_{\text{Sharpe}}
                   + 0.30 \cdot S_{\text{Stabilité}}
                   + 0.20 \cdot S_{\text{Cohérence}}
                   + 0.15 \cdot S_{\text{Drawdown}}
\label{eq:reliability}
\end{equation}

\textbf{Composante 1 — Sharpe (35\,\%)}~:
\begin{equation}
S_{\text{Sharpe}} = \text{clamp}\!\left(\frac{\overline{\text{Sharpe}} + 0.5}{2.0},\; 0,\; 1\right)
\label{eq:sharpe_score}
\end{equation}

Le Sharpe moyen sur $[-0.5, +1.5]$ est linéairement transformé en $[0, 1]$.
Un $\overline{\text{Sharpe}} = 0.5$ correspond à un score de $0.50$ (décent).

\textbf{Composante 2 — Stabilité (30\,\%)}~:
\begin{equation}
S_{\text{Stabilité}} = \frac{\#\{k : \text{Sharpe}_k > 0\}}{n_{\text{fenêtres}}}
\label{eq:stability_score}
\end{equation}

Fraction des sous-périodes avec Sharpe positif. Directement interprétable~:
un signal \og bon \fg{} 70\,\% du temps a un score de 0.70.

\textbf{Composante 3 — Cohérence (20\,\%)}~:
\begin{equation}
S_{\text{Cohérence}} = \text{clamp}\!\left(1 - \frac{\sigma_{\text{Sharpe}}}{|\overline{\text{Sharpe}}| + 0.1},\; 0,\; 1\right)
\label{eq:consistency_score}
\end{equation}

Pénalise les signaux dont le Sharpe varie fortement d'une fenêtre à l'autre.
Le terme $+0.1$ au dénominateur prévient la division par zéro quand
$\overline{\text{Sharpe}} \approx 0$.

\textbf{Composante 4 — Drawdown (15\,\%)}~:
\begin{equation}
S_{\text{Drawdown}} = \text{clamp}\!\left(1 - \frac{\overline{\text{MDD}}}{0.30},\; 0,\; 1\right)
\label{eq:drawdown_score}
\end{equation}

Un drawdown moyen de 30\,\% ou plus donne un score de 0. Le seuil de 30\,\% est
calibré sur la volatilité typique des actions MASI.

\subsection{Justification des poids}

Le Sharpe constitue la mesure principale (35\,\%) car il résume le couple rendement/risque
en une seule valeur. Il n'est toutefois pas suffisant~: un signal avec Sharpe élevé sur
une seule fenêtre et négatif sur les quatre autres est peu fiable. La stabilité (30\,\%)
capture cette dimension temporelle. La cohérence (20\,\%) et le drawdown (15\,\%) sont
des métriques complémentaires qui résistent aux configurations pathologiques.

\begin{quote}
\textit{«~Any metric that becomes a target ceases to be a good metric.~»}
— Goodhart's Law, appliqué au backtesting.
\end{quote}

%% ============================================================
\section{Layers D et E --- Filtrage et Réduction de Redondance}
\label{sec:layers_de}

\subsection{Layer D — Filtre de Survie}

Le filtre de survie est une procédure en deux étapes~:

\begin{enumerate}
    \item \textbf{Porte de viabilité} (Layer~C)~: seules les variantes avec
    $\text{viable} = \text{True}$ passent.

    \item \textbf{Floor absolu}~: seules les variantes avec
    $\text{reliability\_score} \geq 0.25$ sont retenues.

    \textit{Justification}~: éviter le \og best of a bad lot \fg{} — si toutes les
    variantes sont mauvaises, le floor absolu renvoie une liste vide et le signal de
    la famille sera neutre ($= 0$), plutôt que de sélectionner la moins mauvaise.

    \item \textbf{Percentile compétitif}~: parmi les variantes restantes, les 60\,\%
    les mieux classées par reliability\_score sont conservées.
\end{enumerate}

\textbf{Exemple numérique~:}
\begin{itemize}
    \item 30 candidats générés (Layer A)~;
    \item 18 viables (gate Layer C)~;
    \item 15 au-dessus du floor 0.25~;
    \item 9 conservés (top 60\,\% de 15).
\end{itemize}

\subsection{Layer E — Réduction de Redondance}

\begin{quote}
\textit{«~Diversification is the only free lunch in finance.~»}
— Markowitz (1952)~\cite{markowitz1952}, appliqué aux ensembles de signaux.
\end{quote}

Si deux variantes survivantes produisent des tableaux de signaux quasi-identiques,
les inclure toutes les deux dans l'ensemble revient à \emph{compter la même information
deux fois}, ce qui crée une illusion de consensus.

\textbf{Algorithme de clustering glouton par corrélation~:}

\begin{enumerate}
    \item Trier les survivants par reliability\_score (décroissant).
    \item Calculer les tableaux de signaux de toutes les variantes.
    \item Calculer la matrice de corrélations de Pearson par paires.
    \item Sélection gloutonne~:
    \begin{itemize}
        \item Premier représentant~: variante de plus haute fiabilité.
        \item Pour chaque variante suivante~: si $\max_{j \in \text{repres.}} |r_{ij}| \leq 0.85$, accepter~; sinon rejeter.
    \end{itemize}
    \item Arrêt~: 10 représentants sélectionnés ou toutes les variantes traitées.
\end{enumerate}

\begin{figure}[H]
\centering
\includegraphics[width=0.9\linewidth]{generated/ch3_correlation_heatmap.png}
\caption{Principe de réduction de redondance par corrélation~: les variantes trop
corrélées sont remplacées par quelques représentants. Source~: couche~E du Signal
Engine~\cite{repo_signal_docs}.}
\label{fig:correlation_heatmap}
\end{figure}

%% ============================================================
\section{Layers F et G --- Signal Courant et Ensemble}
\label{sec:layers_fg}

\subsection{Layer F — Signal Courant par Représentant}

Pour chaque représentant sélectionné en Layer~E, le signal courant est le signal
calculé sur la \emph{dernière barre disponible}~:

\begin{enumerate}
    \item Calcul du tableau de signal complet~: $\text{sig\_arr} = \texttt{compute\_signal}(\text{close})$~;
    \item Extraction du dernier bar~: $s = \text{sig\_arr}[-1]$~;
    \item Calcul de la valeur brute de l'indicateur (ex. valeur de la SMA, niveau RSI)~;
    \item Assignation du label selon la famille~:
\end{enumerate}

\begin{table}[H]
\centering
\caption{Labels de signal par catégorie selon le score d'ensemble (mêmes seuils que la couche G)}
\label{tab:labels_signal}
\begin{tabularx}{\linewidth}{@{}lYYYYY@{}}
\toprule
\textbf{Score} & \textbf{Tendance} & \textbf{Momentum} & \textbf{Oscillation} & \textbf{Volume} & \textbf{Global} \\
\midrule
$> 50$         & Très haussier     & Fort momentum haussier & Très survendu & Forte accumulation & Achat fort \\
$> 15$         & Haussier          & Momentum haussier      & Survendu      & Accumulation       & Achat \\
$\geq -15$     & Neutre            & Pas de momentum        & Normal        & Neutre             & Neutre \\
$\geq -50$     & Baissier          & Momentum baissier      & Suracheté     & Distribution       & Vente \\
$< -50$        & Très baissier     & Fort momentum baissier & Très suracheté & Forte distribution & Vente forte \\
\bottomrule
\end{tabularx}
\end{table}

\subsection{Layer G — Combinaison Pondérée (Ensemble)}
\label{sec:signaux_ensemble}

Le score d'ensemble est une moyenne pondérée par la fiabilité~:

\begin{equation}
\text{score\_pct} = 100 \times \frac{\displaystyle\sum_{i} s_i \cdot w_i}{\displaystyle\sum_{i} w_i}
\label{eq:ensemble}
\end{equation}

où $s_i \in \{-1, 0, +1\}$ est le signal courant du représentant $i$ et
$w_i = \text{reliability\_score}(i)$ est son poids.

Le score final est dans $[-100, +100]$~:
\begin{itemize}
    \item Proche de $+100$~: consensus fort à l'achat~;
    \item Proche de $-100$~: consensus fort à la vente~;
    \item Proche de $0$~: absence de signal clair.
\end{itemize}

\begin{table}[H]
\centering
\caption{Exemple de calcul du score ensemble pour la famille Tendance sur IAM
(données illustratives)}
\label{tab:ensemble_exemple}
\begin{tabularx}{\linewidth}{@{}lYYYY@{}}
\toprule
\textbf{Variante} & \textbf{Signal $s_i$} & \textbf{Fiabilité $w_i$} & \textbf{Contribution $s_i w_i$} & \textbf{Poids} \\
\midrule
SMA(20)     & $+1$ & 0.72 & $+0.72$ & \\
SMA(50)     & $+1$ & 0.61 & $+0.61$ & \\
SMA(75)     & $-1$ & 0.52 & $-0.52$ & \\
EMA(30)     & $+1$ & 0.65 & $+0.65$ & \\
EMA Cross   &  $0$ & 0.48 & $\phantom{+}0.00$ & \\
\midrule
\textbf{Total} & & $\sum w_i = 2.98$ & $\sum s_i w_i = 1.46$ & \\
\midrule
\multicolumn{4}{l}{\textbf{Score ensemble} $= 100 \times 1.46 / 2.98$} & $\mathbf{= +49.0}$ \\
\bottomrule
\end{tabularx}
\end{table}

Un score de $+49$ indique un biais haussier (label : "Achat", entre 15 et 50), avec un consensus modéré
entre les représentants de la famille Tendance.

%% ============================================================
\section{Limites et Perspectives du Moteur de Signaux}
\label{sec:limites_signaux}

Le pipeline A→G réduit substantiellement le biais de tests multiples mais ne
l'élimine pas totalement. Trois techniques prévues mais non encore intégrées
représentent des améliorations futures~:

\begin{enumerate}
    \item \textbf{Deflated Sharpe Ratio} (Bailey \& López de Prado, 2012~\cite{bailey2012})~:
    ajuste le Sharpe pour le nombre de trials effectués — prévu au Layer~C.
    \item \textbf{Combinatorially Purged Cross-Validation (CPCV)}
    (López de Prado, 2018, Ch.~12)~: élimination de la fuite de données
    temporelle dans la cross-validation.
    \item \textbf{Feature Importance Ranking}
    (López de Prado, 2020, Ch.~6)~: identifier quels indicateurs contribuent
    réellement et lesquels sont redondants à l'échelle du portefeuille.
\end{enumerate}

\section*{Conclusion du chapitre}

Ce chapitre a présenté le premier moteur de résultat de la plateforme~: le Signal Engine,
pipeline de sept couches (A→G) qui transforme 30 candidats par famille en un score
d'ensemble pondéré par sa stabilité temporelle. Chaque couche intègre un mécanisme de défense
spécifique contre l'overfitting~: évaluation systématique par sous-périodes, gate de
viabilité, score multi-composante, floor absolu, clustering de corrélation et moyennage.
Ce moteur alimente directement le tableau de bord (Chapitre~\ref{ch:dashboard}).
Le second moteur, la Walk-Forward Analysis avancée, fait l'objet du chapitre suivant.
